\documentclass[a4paper,12pt]{article}

\usepackage[T2A]{fontenc}
\usepackage[utf8]{inputenc}
\usepackage[russian]{babel}

\usepackage{amsmath,amssymb,mathtools}
\usepackage{helvet}
\renewcommand{\familydefault}{\sfdefault}

\usepackage{icomma}
\usepackage[a4paper,margin=2.2cm,top=2.3cm,bottom=2.3cm]{geometry}
\usepackage{microtype}

\usepackage{tikz}
\usetikzlibrary{calc,positioning}

\usepackage{tabularx,booktabs}
\usepackage{enumitem}
\usepackage{hyperref}
\usepackage{parskip}
\usepackage{fancyhdr}
\usepackage{setspace}
\usepackage{xcolor}

\definecolor{accent}{HTML}{008080}
\definecolor{darkaccent}{HTML}{004D4D}
\definecolor{textgray}{HTML}{333333}
\definecolor{softgray}{HTML}{6B7280}
\definecolor{lightaccent}{HTML}{E8F3F3}

\hypersetup{
    colorlinks=true,
    linkcolor=darkaccent,
    urlcolor=darkaccent
}

\onehalfspacing
\setlength{\parindent}{0pt}
\setlength{\headheight}{15pt}

\pagestyle{fancy}
\fancyhf{}
\fancyhead[L]{\small\color{softgray}Теория вероятностей}
\fancyhead[R]{\small\color{softgray}Ксения Клыкова}
\fancyfoot[C]{\small\color{softgray}\thepage}
\renewcommand{\headrulewidth}{0.3pt}
\renewcommand{\footrulewidth}{0pt}

\setlist[itemize]{
    leftmargin=1.6em,
    itemsep=0.25em,
    topsep=0.35em
}

\setlist[enumerate]{
    leftmargin=2em,
    itemsep=0.5em,
    topsep=0.4em
}

\newcommand{\accentline}{%
    \vspace{0.3em}
    {\color{accent}\rule{\linewidth}{0.7pt}}
    \vspace{0.5em}
}

\newcommand{\important}[1]{\textbf{\color{darkaccent}#1}}

\newcommand{\prob}[1]{P\!\left(#1\right)}

\begin{document}

% ============================================================
% ТИТУЛЬНЫЙ ЛИСТ
% ============================================================

\begin{titlepage}
\thispagestyle{empty}

\begin{center}

\vspace*{2.1cm}

{\Large\color{softgray}Учебный чек-лист}

\vspace{1cm}

{\fontsize{27}{32}\selectfont\bfseries\color{darkaccent}
ТЕОРИЯ ВЕРОЯТНОСТЕЙ}

\vspace{0.45cm}

{\Large
Сложение и умножение вероятностей\\
Дерево вероятностей\\
Схема Бернулли}

\vspace{1.2cm}

{\color{accent}\rule{0.68\linewidth}{1pt}}

\vspace{1.3cm}

{\large
Ключевые идеи, алгоритмы и типовые задачи\\
для системного повторения}

\vfill

{\large
\textbf{Лёвин Артём Александрович}\\[0.25cm]
эксклюзивно для \textbf{Ксении Клыковой}
}

\vspace{0.7cm}

{\large \today}

\vspace{1.8cm}

\begin{tikzpicture}[x=1cm,y=1cm]
    \draw[
        color=accent,
        line width=1.2pt
    ]
    (-6,0)
    .. controls (-3.8,0.75) and (-1.9,-0.55) ..
    (0,0)
    .. controls (2.0,0.55) and (4.0,-0.65) ..
    (6,0);
\end{tikzpicture}

\vspace{0.4cm}

\end{center}
\end{titlepage}

% ============================================================
% ОСНОВНОЙ ЧЕК-ЛИСТ
% ============================================================

\section*{\color{darkaccent}1. Главная идея занятия}
\accentline

Большинство задач этого блока удобно переводить с обычного русского языка
на \important{язык событий}.

Главные смысловые связки:

\begin{center}
\renewcommand{\arraystretch}{1.35}
\begin{tabularx}{0.91\linewidth}{@{}>{\bfseries}l X X@{}}
\toprule
Слово & Смысл & Типичное действие \\
\midrule
«и» &
события должны произойти совместно &
умножаем вероятности вдоль одной последовательности событий \\

«или» &
подходит один из нескольких несовместимых вариантов &
складываем вероятности подходящих вариантов \\

«не» &
происходит противоположное событие &
переходим к дополнению: $1-P(A)$ \\

«хотя бы один» &
один, два, три и т.\,д. успеха &
часто удобнее использовать противоположное событие «ни одного» \\
\bottomrule
\end{tabularx}
\end{center}

\important{Ключевой навык:} прежде чем считать, сформулируйте событие словами
«и», «или», «не», «хотя бы один».

% ------------------------------------------------------------

\section*{\color{darkaccent}2. Противоположные события}
\accentline

Если событие $A$ означает «успех», то противоположное ему событие
$\overline A$ означает «успех не произошёл».

\[
\boxed{\prob{\overline A}=1-\prob{A}}
\]

Например, если вероятность попадания равна $0,8$, то вероятность промаха:

\[
1-0,8=0,2.
\]

Аналогично, если продавец занят с вероятностью $0,1$, то он свободен с
вероятностью

\[
1-0,1=0,9.
\]

\important{Проверка здравого смысла:}
вероятности двух противоположных событий всегда дают единицу:

\[
\prob{A}+\prob{\overline A}=1.
\]

% ------------------------------------------------------------

\section*{\color{darkaccent}3. Союз «и»: умножение вероятностей}
\accentline

Для независимых событий $A$ и $B$:

\[
\boxed{\prob{A\cap B}=\prob{A}\cdot\prob{B}}
\]

Это соответствует формулировке:

\begin{center}
\textbf{«произошло событие $A$ и произошло событие $B$».}
\end{center}

\subsection*{Пример 1. Два биатлониста}

Первый биатлонист попадает с вероятностью $0,2$, второй --- с вероятностью
$0,3$. Выстрелы независимы.

Вероятность того, что попадут \important{оба}:

\[
0,2\cdot0,3=0,06.
\]

Здесь слово «оба» означает:

\[
\text{попал первый \textbf{и} попал второй}.
\]

\subsection*{Пример 2. Заданная последовательность попаданий}

Биатлонист стреляет по пяти мишеням. Вероятность попадания в каждую равна
$0,8$.

Нужно найти вероятность события:

\[
\text{попал в первые три \textbf{и} промахнулся по последним двум}.
\]

Получаем:

\[
0,8^3\cdot0,2^2=0,02048.
\]

Здесь порядок результатов \important{полностью задан}, поэтому число
различных вариантов выбирать не требуется.

% ------------------------------------------------------------

\section*{\color{darkaccent}4. Союз «или»: сложение вероятностей}
\accentline

Если события $A$ и $B$ \important{несовместимы}, то есть одновременно
произойти не могут, то

\[
\boxed{\prob{A\cup B}=\prob{A}+\prob{B}}.
\]

В задачах с несколькими маршрутами логика обычно такая:

\[
\text{маршрут 1 \textbf{или} маршрут 2 \textbf{или} маршрут 3}.
\]

Вероятность каждого отдельного маршрута находится умножением, после чего
вероятности подходящих маршрутов складываются.

\subsection*{Важно}

Формула простого сложения применяется именно для
\important{несовместимых} вариантов.

В общем случае:

\[
\prob{A\cup B}
=
\prob{A}+\prob{B}-\prob{A\cap B}.
\]

% ------------------------------------------------------------

\section*{\color{darkaccent}5. «Хотя бы один»: главный приём через дополнение}
\accentline

Фраза \important{«хотя бы один»} означает:

\[
1,\ 2,\ 3,\ldots
\]

успеха.

Перебирать все подходящие случаи часто неудобно.

Гораздо полезнее правило:

\[
\boxed{
P(\text{хотя бы один успех})
=
1-P(\text{ни одного успеха})
}
\]

или, в удобной словесной форме,

\[
\boxed{
\text{хотя бы один «да»}
=
1-\text{все «нет»}
}.
\]

Аналогично,

\[
\boxed{
\text{хотя бы один «нет»}
=
1-\text{все «да»}
}.
\]

\subsection*{Пример 1. Два биатлониста}

Вероятности попадания:

\[
p_1=0,2,
\qquad
p_2=0,3.
\]

Вероятности промаха:

\[
q_1=0,8,
\qquad
q_2=0,7.
\]

Никто не попал:

\[
0,8\cdot0,7=0,56.
\]

Значит,

\[
P(\text{хотя бы один попал})
=
1-0,56
=
0,44.
\]

Можно проверить ответ полным перебором:

\[
0,2\cdot0,7
+
0,8\cdot0,3
+
0,2\cdot0,3
=
0,44.
\]

\subsection*{Пример 2. Три продавца}

Каждый из трёх продавцов независимо занят с вероятностью $0,1$.

Нужно найти вероятность того, что \important{хотя бы один продавец свободен}.

Противоположное событие: \important{заняты все трое}.

\[
P(\text{все заняты})
=
0,1^3
=
0,001.
\]

Следовательно,

\[
P(\text{хотя бы один свободен})
=
1-0,001
=
0,999.
\]

Приём работает для любого количества независимых событий.

Если вероятности успеха одинаковы и равны $p$, то

\[
\boxed{
P(\text{хотя бы один успех})
=
1-(1-p)^n
}.
\]

% ------------------------------------------------------------

\section*{\color{darkaccent}6. Полная группа событий}
\accentline

Набор несовместимых событий образует \important{полную группу}, если одно
из них обязательно происходит.

Сумма вероятностей событий полной группы равна единице:

\[
\boxed{
P(A_1)+P(A_2)+\ldots+P(A_n)=1
}.
\]

Например, для двух независимых стрелков возможны четыре исхода:

\begin{center}
\renewcommand{\arraystretch}{1.35}
\begin{tabularx}{0.83\linewidth}{@{}X c@{}}
\toprule
Исход & Вероятность \\
\midrule
первый попал, второй попал &
$0,2\cdot0,3=0,06$ \\

первый попал, второй промахнулся &
$0,2\cdot0,7=0,14$ \\

первый промахнулся, второй попал &
$0,8\cdot0,3=0,24$ \\

оба промахнулись &
$0,8\cdot0,7=0,56$ \\
\midrule
Сумма &
$1$ \\
\bottomrule
\end{tabularx}
\end{center}

Такая таблица полезна, когда нужно увидеть все возможные варианты сразу.

% ------------------------------------------------------------

\section*{\color{darkaccent}7. Схема Бернулли}
\accentline

Пусть проводится $n$ независимых испытаний, причём в каждом:

\[
P(\text{успех})=p,
\qquad
P(\text{неудача})=q=1-p.
\]

Вероятность ровно $k$ успехов вычисляется по формуле Бернулли:

\[
\boxed{
P_n(k)
=
C_n^k p^k q^{\,n-k}
}
\]

где

\[
C_n^k
=
\frac{n!}{k!(n-k)!}.
\]

\subsection*{Почему появляется $C_n^k$?}

Выражение

\[
p^kq^{n-k}
\]

задаёт вероятность \important{одного конкретного порядка} из $k$ успехов и
$n-k$ неудач.

Число различных расположений этих $k$ успехов среди $n$ испытаний равно

\[
C_n^k.
\]

Поэтому вероятности всех таких несовместимых последовательностей
складываются.

\subsection*{Пример}

Биатлонист стреляет пять раз, вероятность попадания равна $0,8$.
Найдём вероятность ровно трёх попаданий:

\[
P_5(3)
=
C_5^3
\cdot
0,8^3
\cdot
0,2^2.
\]

Так как

\[
C_5^3=10,
\]

получаем

\[
P_5(3)
=
10\cdot0,8^3\cdot0,2^2
=
0,2048.
\]

\important{Сравните две задачи:}

\[
\begin{aligned}
&\text{«попал именно в первые три»}
&&\longrightarrow
&&0,8^3\cdot0,2^2,
\\[2mm]
&\text{«попал ровно три раза из пяти»}
&&\longrightarrow
&&C_5^3\cdot0,8^3\cdot0,2^2.
\end{aligned}
\]

Это принципиальное различие.

% ------------------------------------------------------------

\section*{\color{darkaccent}8. Дерево вероятностей}
\accentline

Дерево вероятностей --- способ представить последовательность случайных
выборов в виде развилок.

Основные правила:

\begin{enumerate}
    \item На каждой развилке выпишите все возможные варианты.
    \item Подпишите вероятность каждой ветви.
    \item Вероятности ветвей, выходящих из одного узла, должны давать $1$.
    \item Вероятность одной полной траектории находится
    \important{умножением} вероятностей вдоль неё.
    \item Если подходят несколько взаимоисключающих траекторий, их
    вероятности \important{складываются}.
\end{enumerate}

\begin{center}
\begin{tikzpicture}[
    >=latex,
    every node/.style={font=\small},
    dot/.style={
        circle,
        fill=accent,
        inner sep=1.6pt
    },
    terminal/.style={
        draw=accent,
        rounded corners=2pt,
        inner sep=4pt
    }
]

\node[dot] (S) at (0,0) {};
\node[left=4pt of S] {$S$};

\node[dot] (A) at (3,1.4) {};
\node[dot] (B) at (3,-1.4) {};

\node[terminal] (AG) at (6,2.2) {успех};
\node[terminal] (AB) at (6,0.7) {неудача};
\node[terminal] (BG) at (6,-0.7) {успех};
\node[terminal] (BB) at (6,-2.2) {неудача};

\draw[->,line width=0.8pt]
    (S)--node[above,sloped] {$p_1$}(A);

\draw[->,line width=0.8pt]
    (S)--node[below,sloped] {$1-p_1$}(B);

\draw[->,line width=0.8pt]
    (A)--node[above,sloped] {$p_2$}(AG);

\draw[->,line width=0.8pt]
    (A)--node[below,sloped] {$1-p_2$}(AB);

\draw[->,line width=0.8pt]
    (B)--node[above,sloped] {$p_2$}(BG);

\draw[->,line width=0.8pt]
    (B)--node[below,sloped] {$1-p_2$}(BB);

\end{tikzpicture}
\end{center}

Например, вероятность пройти по верхней траектории к успеху:

\[
p_1\cdot p_2.
\]

Если к нужному результату ведут несколько путей:

\[
P
=
P(\text{путь 1})
+
P(\text{путь 2})
+
\ldots
\]

% ------------------------------------------------------------

\section*{\color{darkaccent}9. Равновероятный выбор на развилках}
\accentline

Если из некоторой точки можно продолжить движение по $m$ равноправным
дорогам и выбор совершается случайно, вероятность выбрать каждую дорогу:

\[
\frac{1}{m}.
\]

Например:

\begin{itemize}
    \item из двух дорог нужная одна:
    \[
    P=\frac12;
    \]

    \item из трёх дорог нужная одна:
    \[
    P=\frac13;
    \]

    \item для траектории с последовательными вероятностями
    $\frac12$, $\frac13$, $\frac12$:
    \[
    P
    =
    \frac12\cdot\frac13\cdot\frac12
    =
    \frac1{12}.
    \]
\end{itemize}

При движении по одной траектории используется союз \important{«и»}, поэтому
вероятности перемножаются.

% ------------------------------------------------------------

\section*{\color{darkaccent}10. Дерево для задач о заводах, товарах и качестве}
\accentline

Типовая конструкция:

\[
\text{источник товара}
\quad\longrightarrow\quad
\text{качество товара}.
\]

Рассмотрим пример.

Первый завод выпускает $40\%$ ламп, второй --- $60\%$.

Среди продукции первого завода $2\%$ брака, среди продукции второго ---
$3\%$.

Тогда:

\[
\begin{aligned}
P(Z_1)&=0,4,
&
P(Z_2)&=0,6,
\\
P(Q\mid Z_1)&=0,98,
&
P(Q\mid Z_2)&=0,97.
\end{aligned}
\]

Событие «лампа исправна» возможно двумя способами:

\[
\text{первый завод и качество}
\]

или

\[
\text{второй завод и качество}.
\]

Поэтому

\[
\begin{aligned}
P(Q)
&=
0,4\cdot0,98
+
0,6\cdot0,97
\\
&=
0,392+0,582
\\
&=
0,974.
\end{aligned}
\]

\[
\boxed{P(Q)=0,974}
\]

То есть вероятность купить исправную лампу равна $97,4\%$.

Это частный случай \important{формулы полной вероятности}:

\[
\boxed{
P(A)
=
P(H_1)P(A\mid H_1)
+
P(H_2)P(A\mid H_2)
+\ldots
}
\]

где $H_1,H_2,\ldots$ --- несовместимые случаи, образующие полную группу.

% ------------------------------------------------------------

\section*{\color{darkaccent}11. Универсальный алгоритм}
\accentline

Получив задачу, действуйте в таком порядке.

\begin{enumerate}
    \item \important{Определите случайные события.}

    Что именно может произойти?

    \item \important{Переведите условие на язык союзов.}

    Где встречаются «и», «или», «не», «хотя бы один»?

    \item \important{Найдите вероятности противоположных событий.}

    Если известно $p$, то
    \[
    q=1-p.
    \]

    \item \important{Проверьте независимость.}

    Простое перемножение вероятностей нескольких событий допустимо, если
    условие позволяет считать их независимыми либо если на дереве уже
    указаны соответствующие условные вероятности.

    \item \important{Разбейте задачу на траектории.}

    Одна траектория:
    \[
    \text{умножение}.
    \]

    Несколько подходящих взаимоисключающих траекторий:
    \[
    \text{сложение}.
    \]

    \item \important{Ищите противоположное событие.}

    Особенно при словах:
    \[
    \text{«хотя бы один»},\quad
    \text{«не все»},\quad
    \text{«есть хотя бы один»}.
    \]

    \item \important{Если требуется ровно $k$ успехов из $n$, проверьте
    возможность схемы Бернулли.}

    \[
    P_n(k)=C_n^kp^k(1-p)^{n-k}.
    \]

    \item \important{Проверьте ответ.}

    Вероятность обязана удовлетворять
    \[
    0\le P\le1.
    \]
\end{enumerate}

% ------------------------------------------------------------

\section*{\color{darkaccent}12. Типичные ошибки}
\accentline

\begin{itemize}

    \item \important{Механически складывать вероятности при слове «или».}

    Сначала нужно проверить, могут ли события произойти одновременно.

    \item \important{Механически умножать вероятности при слове «и».}

    Для независимых событий это допустимо. В дереве нужно использовать
    вероятность именно той ветви, которая выходит из текущего узла.

    \item \important{Забывать перейти от успеха к неудаче.}

    Если попадание имеет вероятность $0,8$, промах имеет вероятность
    $0,2$.

    \item \important{Путать «ровно три» и «первые три».}

    \[
    \text{первые три успеха}
    \neq
    \text{ровно три успеха из пяти}.
    \]

    Во втором случае учитываются разные расположения успехов.

    \item \important{Перебирать все случаи для «хотя бы одного».}

    Обычно быстрее:
    \[
    1-P(\text{ни одного}).
    \]

    \item \important{Складывать обыкновенные дроби по числителям и
    знаменателям.}

    Например,
    \[
    \frac1{16}+\frac1{24}
    \neq
    \frac2{40}.
    \]

    Нужно привести дроби к общему знаменателю:

    \[
    \frac1{16}+\frac1{24}
    =
    \frac3{48}+\frac2{48}
    =
    \frac5{48}.
    \]

    \item \important{Принимать дерево вероятностей за отдельную формулу.}

    Дерево --- это способ организации информации. Расчёты внутри него
    основаны на тех же правилах сложения и умножения вероятностей.

\end{itemize}

% ------------------------------------------------------------

\section*{\color{darkaccent}13. Мини-чек перед завершением задачи}
\accentline

Перед записью ответа задайте себе пять вопросов:

\begin{enumerate}
    \item Что именно считается успехом?
    \item Есть ли здесь событие, противоположное искомому?
    \item Где в условии стоит логическое «и»?
    \item Где стоит логическое «или»?
    \item Нужно ли учитывать несколько различных расположений одинакового
    количества успехов?
\end{enumerate}

Если ответы на эти вопросы ясны, большая часть расчёта уже определена.

% ============================================================
% ТРЕНИРОВОЧНЫЙ БЛОК
% ============================================================

\clearpage

\section*{\color{darkaccent}Тренировочный блок}
\accentline

\textbf{Ксения, цель блока:}
научиться быстро переводить условие задачи на язык событий и выбирать
между умножением, сложением, противоположным событием, деревом
вероятностей и схемой Бернулли.

Решения в этом разделе записывать не требуется. Рекомендуется рядом с
каждым выражением кратко подписывать используемый принцип:
\textit{«и»}, \textit{«или»}, \textit{«дополнение»},
\textit{«Бернулли»}.

\subsection*{\color{darkaccent}Средний уровень}

\begin{enumerate}

\item
Первый стрелок поражает мишень с вероятностью $0,7$, второй --- с
вероятностью $0,6$. Стрелки действуют независимо.

Найдите вероятность того, что мишень поразят оба стрелка.

\item
Вероятность того, что каждый из четырёх независимо работающих датчиков
исправен, равна $0,9$.

Найдите вероятность того, что хотя бы один датчик окажется неисправен.

\item
Спортсмен выполняет пять независимых попыток. Вероятность успеха в каждой
равна $0,75$.

Найдите вероятность того, что первые три попытки будут успешными, а
последние две --- неуспешными.

\end{enumerate}

\subsection*{\color{darkaccent}Выше среднего}

\begin{enumerate}[resume]

\item
Биатлонист стреляет по пяти мишеням. Вероятность попадания при каждом
выстреле равна $0,8$.

Найдите вероятность ровно трёх попаданий.

\item
Три игрока независимо выполняют по одной попытке. Вероятности успеха равны

\[
0,4,\qquad 0,5,\qquad 0,7.
\]

Найдите вероятность того, что хотя бы один игрок добьётся успеха.

\item
На первой развилке путешественник выбирает одну из двух дорог
равновероятно. Если он выбирает нужную дорогу, то на следующей развилке
ему необходимо выбрать одну определённую дорогу из трёх равновероятных.

Найдите вероятность пройти по требуемому маршруту.

\item
К нужной точке ведут два несовместимых маршрута.

Первый маршрут требует последовательно выбрать нужные ветви с
вероятностями

\[
\frac12,\qquad \frac13,\qquad \frac12.
\]

Второй маршрут требует выборов с вероятностями

\[
\frac12,\qquad \frac14.
\]

Найдите вероятность попасть в нужную точку.

\item
Первое предприятие выпускает $35\%$ деталей, второе --- $65\%$.
Доля брака на первом предприятии составляет $4\%$, на втором --- $2\%$.

Найдите вероятность того, что случайно выбранная деталь окажется
исправной.

\item
Вероятность правильного ответа на каждый из шести независимых вопросов
равна $0,7$.

Найдите вероятность того, что ученик правильно ответит ровно на четыре
вопроса.

\end{enumerate}

\subsection*{\color{darkaccent}Сложный уровень}

\begin{enumerate}[resume]

\item
На склад поступают изделия с трёх фабрик:

\[
P(F_1)=0,2,\qquad
P(F_2)=0,5,\qquad
P(F_3)=0,3.
\]

Вероятность того, что изделие окажется качественным, составляет

\[
P(Q\mid F_1)=0,96,
\qquad
P(Q\mid F_2)=0,94,
\qquad
P(Q\mid F_3)=0,98.
\]

После выбора изделия независимо проводится электронная проверка.
Если изделие качественное, проверка успешно завершается с вероятностью
$0,99$.

Найдите вероятность того, что случайно выбранное изделие окажется
качественным \textbf{и} успешно пройдёт электронную проверку.

\end{enumerate}

% ============================================================
% ОТВЕТЫ
% ============================================================

\clearpage

\section*{\color{darkaccent}Ответы к тренировочному блоку}
\accentline

\renewcommand{\arraystretch}{1.45}

\begin{table}[ht]
\centering
\caption{Ответы}
\begin{tabularx}{\linewidth}{@{}c X@{}}
\toprule
\textbf{№} & \textbf{Ответ} \\
\midrule

1 &
\[
0,7\cdot0,6=0,42
\]
\\

2 &
\[
1-0,9^4=0,3439
\]
\\

3 &
\[
0,75^3\cdot0,25^2
=
\frac{27}{1024}
\approx0,02637
\]
\\

4 &
\[
C_5^3\cdot0,8^3\cdot0,2^2
=
0,2048
\]
\\

5 &
\[
1-(1-0,4)(1-0,5)(1-0,7)
=
1-0,6\cdot0,5\cdot0,3
=
0,91
\]
\\

6 &
\[
\frac12\cdot\frac13
=
\frac16
\]
\\

7 &
\[
\frac12\cdot\frac13\cdot\frac12
+
\frac12\cdot\frac14
=
\frac1{12}+\frac18
=
\frac5{24}
\]
\\

8 &
\[
0,35\cdot0,96
+
0,65\cdot0,98
=
0,973
\]
\\

9 &
\[
C_6^4\cdot0,7^4\cdot0,3^2
=
15\cdot0,7^4\cdot0,3^2
=
0,324135
\]
\\

10 &
\[
\begin{aligned}
P(Q)
&=
0,2\cdot0,96
+
0,5\cdot0,94
+
0,3\cdot0,98
\\
&=
0,956,
\\[1mm]
P(Q\cap T)
&=
0,956\cdot0,99
=
0,94644.
\end{aligned}
\]
\\

\bottomrule
\end{tabularx}
\end{table}

\vfill

\begin{center}
{\color{accent}\rule{0.65\linewidth}{0.7pt}}

\vspace{0.5em}

{\small\color{softgray}
Главная идея: одна траектория --- умножение;\\
несколько подходящих несовместимых траекторий --- сложение;\\
«хотя бы один» --- сначала проверьте противоположное событие.
}
\end{center}

\end{document}